\subsection{申请人简历}
\XJField{applicant}，\XJField{sex}，\XJField{birth}出生，\XJField{degree}，\XJField{organization}\XJField{professional-title}、硕士研究生导师。主要从事时间序列分析、新能源功率预测、多尺度时频表示与多源信息融合研究。现所在院系为\XJField{department}。通讯地址：\XJField{address}；邮编：\XJField{postcode}；电子邮箱：\XJField{email}。
% 院系名称、邮编、邮箱及以下履历年月沿用共用档案的公开来源待终核状态，不等同于申请人已终核；详见docs/填写说明与待补清单.md。

\subsubsection{教育经历}
2019.09--2023.08，新疆大学，计算机科学与技术专业，博士。

2012.09--2014.06，上海海事大学，计算机技术专业，硕士。

2008.09--2012.06，海南热带海洋学院，生态学专业，学士。

\subsubsection{工作经历}
2023.12至今，新疆大学，副教授。

2023.10--2023.12，新疆大学，教学与科研工作。

2018.03--2019.08，浙江大学山东工业技术研究院，软件工程师。

2016.10--2017.12，上海中科软科技有限公司，软件工程师。

\subsubsection{与本项目直接相关的代表成果}
申请人作为通讯作者发表DWT-Former研究（\textit{Energy}, 2025, 341: 139283）\cite{sun2025dwt}，作为共同第一作者发表WaveKAN研究（\textit{Neurocomputing}, 2026, 682: 133361）\cite{liu2026wavekan}。上述成果分别支撑光伏多尺度预测和自适应小波长时序建模的方法基础；本项目专属概率校准、爬坡预警及新疆场站先导试验尚无可核验证据。
